\documentclass[11pt,a4paper]{article}
\usepackage[T1]{fontenc}
\usepackage{lmodern,microtype,amsmath,amssymb,amsthm,mathtools}
\usepackage[margin=27mm]{geometry}
\usepackage{booktabs,array,enumitem}
\usepackage[colorlinks=true,linkcolor=blue!50!black,citecolor=blue!50!black,urlcolor=blue!50!black]{hyperref}
\usepackage{xcolor,fancyhdr}
\pagestyle{fancy}\fancyhf{}
\fancyhead[L]{\small Polynomial rigidity for constant terms}
\fancyhead[R]{\small Research note}\fancyfoot[C]{\thepage}
\setlength{\headheight}{14pt}
\setlength{\parskip}{4pt}
\newtheorem{theorem}{Theorem}[section]
\newtheorem{lemma}[theorem]{Lemma}
\newtheorem{corollary}[theorem]{Corollary}
\newtheorem{proposition}[theorem]{Proposition}
\theoremstyle{remark}\newtheorem{remark}[theorem]{Remark}
\newcommand{\Q}{\mathbb Q}\newcommand{\Z}{\mathbb Z}
\newcommand{\CT}{\operatorname{CT}}
\newcommand{\Newt}{\operatorname{Newt}}
\newcommand{\supp}{\operatorname{supp}}
\newcommand{\lcm}{\operatorname{lcm}}
\title{Polynomial Rigidity for Congruence-Preserving\\Constant-Term Sequences}
\author{Research note prepared for KEIO}
\date{November 2026}
\hypersetup{pdftitle={Polynomial Rigidity for Congruence-Preserving Constant-Term Sequences},pdfauthor={Research note prepared for KEIO}}
\begin{document}
\maketitle
\begin{abstract}
Let an integer sequence be a finite rational linear combination of sequences
$\CT(GF^n)$, where $F,G$ are multivariate Laurent polynomials over $\Q$.
We prove that if $a_{n+p}\equiv a_n\pmod p$ for every $n\ge0$ and infinitely
many primes $p$, then $a_n$ is a rational polynomial in $n$.
The result permits polynomial weights in $n$, and has no growth hypothesis.
The proof combines Frobenius digit factorization, finite polynomial
interpolation, and the Duistermaat--van der Kallen constant-term theorem.
It gives a restricted case of Ruzsa's 1971 conjecture, not its general solution.
We also prove sharp limitations of separately uniform arithmetic and analytic
Hankel estimates. The arguments have been internally cross-checked; external
review and priority have not been established.
\end{abstract}

\section{The old problem and the result proved here}

Ruzsa's conjecture \cite{Ruzsa,Delaygue} asserts that an integer sequence
satisfying
\[
 a_{n+k}\equiv a_n\pmod k\quad(n\ge0,\ k\ge1),\qquad
 \limsup_{n\to\infty}|a_n|^{1/n}<e
\]
must have the form $a_n=P(n)$ with $P\in\Q[n]$.
Its 1971 origin meets the twenty-year age requirement. Delaygue's 2026 paper
still treats the general case as open and proves the case of at most two
singular directions of the generating function. The corresponding
\emph{primary} congruences require the displayed relation only for prime $k$.

For a finite Laurent polynomial $A$ in any number of variables, $\CT(A)$
denotes the coefficient of the monomial with every exponent zero.
Different summands below may use different sets of variables.

\begin{theorem}[Polynomial rigidity]\label{thm:main}
Suppose $a_n\in\Z$ for all $n\ge0$, and
\begin{equation}\label{eq:representation}
 a_n=\sum_{j=1}^s q_j(n)\CT\bigl(G_j F_j^n\bigr),
 \qquad q_j\in\Q[n],\quad
 F_j,G_j\in\Q[x_1^{\pm1},\ldots,x_{d_j}^{\pm1}].
\end{equation}
Assume that there is an infinite set of primes $\mathcal P$ such that
\begin{equation}\label{eq:period}
 a_{n+p}\equiv a_n\pmod p
 \quad\text{for every }n\ge0\text{ and every }p\in\mathcal P.
\end{equation}
Then $a_n=P(n)$ for a polynomial $P\in\Q[n]$, for every $n\ge0$.
No bound on the growth of $a_n$ is required.
\end{theorem}

More explicitly, define
\[
 b_j(m)=\CT(F_j^m),\qquad
 J=\{j:b_j(m)=1\text{ for every }m\ge0\}.
\]
The proof gives the exact projection
\begin{equation}\label{eq:projection}
 a_n=\sum_{j\in J}q_j(n)\CT(G_jF_j^n).
\end{equation}
Every summand on the right is polynomial in $n$.
The condition $j\in J$ can equivalently be expressed as
$F_j=1$ or $0\notin\Newt(F_j-1)$.

\begin{corollary}[Finite or cofinite prime periods]\label{cor:finitelymany}
For a nonpolynomial integer sequence of the form
\eqref{eq:representation}, only finitely many primes $p$ can satisfy
$a_{n+p}\equiv a_n\pmod p$ for all $n\ge0$.
For a polynomial sequence the same property holds for all but finitely many
primes. Thus the set of such primes is always finite or cofinite in this class.
In particular, Ruzsa's conclusion holds throughout this class, even after
its growth hypothesis is removed.
\end{corollary}

\begin{proof}
The nonpolynomial assertion is the contrapositive of Theorem~\ref{thm:main}.
For a rational polynomial, reduction modulo any prime avoiding its
coefficient denominators gives $P(n+p)=P(n)$ in $\mathbb F_p$.
\end{proof}

The quantifier ``for all $n$'' is essential. A congruence at one index,
or a finite numerical check, is not the hypothesis of the theorem.

\section{Frobenius digits and finite interpolation}

\begin{lemma}[Digit factorization]\label{lem:digits}
Fix Laurent polynomials $F,G$ over $\Q$ and fixed nonnegative integers
$r,m_1,\ldots,m_k$. For all sufficiently large primes $p$, put
$N_p=r+\sum_{i=1}^k m_i p^i$. Then
\begin{equation}\label{eq:digit}
 \CT(GF^{N_p})\equiv
 \CT(GF^r)\prod_{i=1}^k\CT(F^{m_i})\pmod p.
\end{equation}
Rational coefficients are reduced only at primes avoiding their denominators.
\end{lemma}

\begin{proof}
For such primes, Frobenius gives
\[
 GF^{N_p}\equiv GF^r\prod_{i=1}^k F(x_1^{p^i},\ldots,x_d^{p^i})^{m_i}
 \pmod p.
\]
Choose $p$ larger than the absolute value of every coordinate of every
exponent appearing in $GF^r$ and in the finitely many $F^{m_i}$.
A contribution to the constant term has exponent vectors satisfying
$v_0+p v_1+\cdots+p^k v_k=0$.
Reduction modulo $p$ and the bound $|(v_0)_h|<p$ imply $v_0=0$.
Divide by $p$ and repeat to obtain $v_i=0$ for every $i$.
Thus constant terms factor. Negative Laurent exponents cause no difficulty.
This proves \eqref{eq:digit}.
\end{proof}

This is an iterated form of the partial Lucas congruences for constant-term
sequences; compare \cite[Lemma 3.1]{BSY}. Its proof above is included to make
the needed hypotheses explicit.

\begin{lemma}[Mixed identities]\label{lem:mixed}
Under the hypotheses of Theorem~\ref{thm:main}, fix $r\ge0$ and set
$c_j(r)=q_j(r)\CT(G_jF_j^r)$. For every finite list
$m_1,\ldots,m_k\ge0$, including the empty list,
\begin{equation}\label{eq:mixed}
 \sum_{j=1}^s c_j(r)\prod_{i=1}^k b_j(m_i)=a_r.
\end{equation}
\end{lemma}

\begin{proof}
Choose primes in the infinite set $\mathcal P$ large enough for
Lemma~\ref{lem:digits}, and avoiding the fixed coefficient denominators.
Since $N_p\equiv r\pmod p$, the global periodicity assumption gives
$a_{N_p}\equiv a_r\pmod p$, while $q_j(N_p)\equiv q_j(r)\pmod p$.
The digit lemma yields \eqref{eq:mixed} modulo every such prime.
Both sides of \eqref{eq:mixed} are fixed rational numbers, independent of $p$.
After clearing a fixed denominator, their difference is an integer divisible
by infinitely many primes. It is therefore zero. The empty-list case is
\eqref{eq:representation} at $n=r$.
\end{proof}

\begin{proposition}[Elimination of nonconstant components]\label{prop:elim}
The mixed identities imply \eqref{eq:projection} for every $n\ge0$.
\end{proposition}

\begin{proof}
If every index lies in $J$, there is nothing to prove. Otherwise enumerate
the remaining indices as $j_1,\ldots,j_t$ and, for each $\ell$, choose
$u_\ell\ge0$ with $b_{j_\ell}(u_\ell)\ne1$. For every $j$ write
\[
 x_j=(b_j(u_1),\ldots,b_j(u_t))\in\Q^t.
\]
Repeated digits in \eqref{eq:mixed}, followed by linearity, show that for
every polynomial $R\in\Q[X_1,\ldots,X_t]$,
\begin{equation}\label{eq:functional}
 \sum_{j=1}^s c_j(r)R(x_j)=a_r R(1,\ldots,1).
\end{equation}
Take the explicit polynomial
\[
 R(X)=\prod_{\ell=1}^t
 \frac{X_\ell-b_{j_\ell}(u_\ell)}{1-b_{j_\ell}(u_\ell)}.
\]
It equals $1$ at $(1,\ldots,1)$, equals $1$ at $x_j$ for $j\in J$, and
vanishes at each $x_{j_\ell}$ because its $\ell$-th factor vanishes.
Equation~\eqref{eq:functional} gives $a_r=\sum_{j\in J}c_j(r)$.
The argument allows arbitrary signs, repeated components, and cancellation
between summands. Since $r$ was arbitrary, the projection holds at every index.
\end{proof}

\section{Why the surviving components are polynomial}

We use the following established theorem, including its geometric conclusion.

\begin{theorem}[Duistermaat--van der Kallen, 1998]\label{thm:dvk}
If a nonzero complex Laurent polynomial $H$ satisfies
$\CT(H^m)=0$ for every integer $m\ge1$, then
$0\notin\Newt(H)$, the convex hull of its exponent support.
Consequently, for every fixed Laurent polynomial $G$,
$\CT(GH^m)=0$ for all sufficiently large $m$.
\end{theorem}

The geometric statement is proved in \cite{DvK}. To see its displayed
consequence, choose a real linear functional $\ell$ and $\eta>0$ with
$\ell(v)\ge\eta$ on $\supp(H)$. Let
$\gamma=\min\{\ell(w):w\in\supp(G)\}$ when $G\ne0$.
Every exponent in $GH^m$ has $\ell$-value at least $\gamma+m\eta$;
when this is positive it cannot be the zero exponent.

\begin{proof}[Completion of Theorem~\ref{thm:main}]
For $j\in J$, put $H_j=F_j-1$. Binomial expansion gives, for $m\ge1$,
\[
 \CT(H_j^m)=\sum_{h=0}^m(-1)^{m-h}\binom mh\CT(F_j^h)
 =\sum_{h=0}^m(-1)^{m-h}\binom mh=0.
\]
If $H_j\ne0$, Theorem~\ref{thm:dvk} supplies an integer $M_j\ge0$ such that
$\CT(G_jH_j^h)=0$ for $h>M_j$. The same conclusion is immediate when
$H_j=0$ or $G_j=0$. Hence for every nonnegative integer $n$,
\begin{equation}\label{eq:finitepoly}
 \CT(G_jF_j^n)=
 \sum_{h=0}^{M_j}\binom nh\CT(G_jH_j^h).
\end{equation}
Here $\binom nh$ is its usual polynomial in $n$; it vanishes for integers
$0\le n<h$, so the formula holds also at the small indices.
Insert \eqref{eq:finitepoly} into \eqref{eq:projection}. This is the required
rational polynomial $P(n)$, valid for all $n\ge0$.
\end{proof}

Conversely, $0\notin\Newt(F-1)$ implies that every positive power of $F-1$
has zero constant term, and therefore $\CT(F^m)=1$ for all $m$.
This justifies the finite geometric test stated after the main theorem.

\begin{corollary}[Pure powers]
If $G_j=1$ for every $j$, then the conclusion is the especially explicit
formula $a_n=\sum_{j\in J}q_j(n)$.
For a single sequence $a_n=\CT(F^n)$, the hypothesis forces $a_n=1$ for all $n$.
\end{corollary}

\section{Scope, examples, and the remaining problem}

The class in Theorem~\ref{thm:main} is closed under finite sums, termwise
products, nonnegative index shifts, and multiplication by a rational
polynomial in $n$. For products, use disjoint variable sets. For a shift by
$h$, replace $G$ by $GF^h$ and $q(n)$ by $q(n+h)$.
Polynomial weights can themselves be absorbed
into constant terms: expand $q(n)=\sum_k d_k\binom nk$ and use
\[
 \binom nk\CT_x(GF^n)
 =\CT_{x,u}\bigl(G u^{-k}[F(x)(1+u)]^n\bigr).
\]
Thus the formulation is equivalent to a finite rational linear combination
of standard constant-term sequences.

For example, the Franel sequence has the representation
\[
 \sum_{k=0}^n\binom nk^3
 =\CT_{x,y}\!\left[\bigl((1+x)(1+y)(1+(xy)^{-1})\bigr)^n\right].
\]
The theorem applies to finite polynomial-weighted combinations of such
sequences and their fixed constant-term weights. It does not assert that
the Franel sequence itself has primary periodicity; the theorem restricts
which combinations can have it.

A simple surviving weighted component is
\[
 F(x)=1+x,\quad G(x)=x^{-h},\qquad
 \CT(GF^n)=\binom nh.
\]
Its Newton polytope is strictly separated from zero after replacing $F$ by
$F-1$. The conclusion allows nonconstant polynomials. The converse
``every integer-valued polynomial is primary'' is not being asserted.

\paragraph{Relation to the corpus.}
The proof develops the constant-term and Frobenius mechanisms used in the
corpus's Franel and lifting work. Its required arithmetic ingredient is
proved in Lemma~\ref{lem:digits}; the substantive external theorem is
Duistermaat--van der Kallen. No unverified assertion from a corpus draft is
needed as a premise. Finite interpolation is the step that permits several
different constant-term components to interact without assuming they are
linearly independent.

\paragraph{Exact limit of the result.}
An arbitrary sequence in Ruzsa's conjecture is not known to have
\eqref{eq:representation}. Even all rational diagonals do not belong to this
class: Bostan--Straub--Yurkevich prove that the Fibonacci sequence is not a
finite rational linear combination of standard constant-term sequences
\cite{BSY}. This is a representation limitation,
not a counterexample to Ruzsa's conjecture. A general proof still needs a
different bridge from congruences and small growth to polynomiality.

\paragraph{Status.}
The main theorem is a complete proof for its stated class. It neither resolves
the full 1971 conjecture nor establishes that this restricted statement was
previously open. Targeted literature searches and internal independent
proof audits are not substitutes for an external priority check or peer review.

\appendix
\section{A sharp arithmetic barrier for Hankel determinants}

Write $H_N(a)=(a_{i+j})_{0\le i,j<N}$, and
$D_N=\prod_{p<N}p^{N-p}$. Let $B_N$ be the positive gcd of all nonzero
$\det H_N(a)$ over integer primary congruence-preserving sequences. Let $C_N$
be the analogous gcd for full congruence preservation.

\begin{theorem}\label{thm:gcd}
For $N\ge2$,
\[
 B_N=D_N,\qquad \log C_N=\tfrac12N^2+o(N^2).
\]
Thus full prime-power congruences cannot improve the quadratic exponent in a
universal divisibility bound obtained from congruences alone.
\end{theorem}

\begin{proof}
Put $b_j=\Delta^j a_0$, so $a_m=\sum_j\binom mj b_j$.
If $T=(\binom ij)_{0\le i,j<N}$, Vandermonde's identity gives
$H_N(a)=T H_N(b)T^{\mathsf T}$, hence the determinants agree.
Let $P_j=\prod_{p\le j}p$ and $L_j=\lcm(1,\ldots,j)$, with $P_0=L_0=1$.
Primary congruences imply $P_j\mid b_j$: modulo $p$, the shift satisfies
$E^p a=a$, so $\Delta^p a=(E-1)^p a=0$. Each row $i$ of $H_N(b)$ is
divisible by $P_i$, giving $D_N\mid\det H_N(a)$.

For a fixed prime $p<N$, choose just
$b_{p-1}=P_{p-1}$ and $b_{N+p-1}=P_{N+p-1}$ nonzero.
The matrix has reverse-diagonal blocks of sizes $p$ and $N-p$, with zero
off-diagonal blocks. Its determinant has $p$-valuation exactly $N-p$.
For $p\ge N$, use only $b_{N-1}=P_{N-1}$, obtaining a $p$-adic unit.
These finite Newton sequences are primary: coefficients with index at least
$p$ vanish modulo $p$, and $\binom mj$ for $j<p$ is polynomial modulo $p$.
This proves $B_N=D_N$.

For full congruences, every sequence with $L_j\mid b_j$ is admissible,
because $L_j\binom mj$ preserves congruences. Indeed, in Vandermonde's
formula for $\binom{m+k}{j}-\binom mj$, each term contains
$L_j\binom ki=k(L_j/i)\binom{k-1}{i-1}$ for some $1\le i\le j$.
The single choice $b_{N-1}=L_{N-1}$ proves
$D_N\mid C_N\mid L_{N-1}^N$.
When $p<N$ and $p^2>2N-2$, the preceding two-block construction with $L$
in place of $P$ still gives $v_p(C_N)=N-p$.
Only primes $p\le\sqrt{2N-2}$ can contribute extra valuations, and
$v_p(C_N)\le N\lfloor\log_p(N-1)\rfloor$. Consequently
\[
 0\le\log C_N-\log D_N
 \le N\pi(\sqrt{2N-2})\log(N-1)=o(N^2).
\]
Finally, the prime number theorem gives
$\log D_N=\sum_{j=0}^{N-1}\sum_{p\le j}\log p=\tfrac12N^2+o(N^2)$.
\end{proof}

The extremizing polynomials depend on $N$ and $p$. This theorem does not
exclude stronger estimates that combine the growth and differential
structure of one fixed nonpolynomial sequence with its congruences.

\section{An exact three-direction analytic barrier}

Consider
\[
 F(z)=(1-20z^3)^{-1/2}=\sum_{n\ge0}a_nz^n,
 \qquad a_{3k}=5^k\binom{2k}{k},\quad a_{3k+1}=a_{3k+2}=0.
\]
This integer, algebraic series has three equally spaced singular directions
and coefficient growth $20^{1/3}<e$. Nevertheless its Hankel determinants
$\mathcal H_N=\det H_N(a)$ satisfy
\begin{align*}
 \mathcal H_{3m}&=(-1)^m2^{3m-1}5^{3m^2-m} &&(m\ge1),\\
 \mathcal H_{3m+1}&=(-1)^m2^{3m}5^{3m^2+m} &&(m\ge0),\\
 \mathcal H_{3m+2}&=0 &&(m\ge0).
\end{align*}
In particular $\limsup_N|\mathcal H_N|^{1/N^2}=5^{1/3}>\sqrt e$.

Here is a proof of the formulas. The numbers $c_k=\binom{2k}{k}$ are the
moments of $dx/(\pi\sqrt{x(4-x)})$ on $(0,4)$.
The monic orthogonal polynomials are $p_0=1$ and
$p_j(x)=2T_j((x-2)/2)$ for $j\ge1$, with squared norms $1,2,2,\ldots$.
Thus, for $m\ge1$, $\det(c_{i+j})_{i,j<m}=2^{m-1}$.
The standard determinant expression for a monic orthogonal polynomial,
evaluated at zero, gives
\[
 p_m(0)=(-1)^m\frac{\det(c_{i+j+1})_{i,j<m}}{\det(c_{i+j})_{i,j<m}}
 =2(-1)^m,
\]
so the shifted determinant is $2^m$. After replacing $c_k$ by $5^kc_k$,
these determinants become $2^{m-1}5^{m(m-1)}$ and $2^m5^{m^2}$.
Reorder the original Hankel rows and columns by their indices modulo three.
The zero-residue block is the first type; the other two residue classes
couple by the second type, with zero diagonal blocks. For $N=3m$ their sizes
are $m,m,m$; for $N=3m+1$, they are $m+1,m,m$; for $N=3m+2$ the last two
sizes differ and force a zero determinant. This proves the claims, with
$N=1,2$ checked directly.

This example is not a Ruzsa counterexample: $a_2=0\not\equiv a_0=1\pmod2$.
It shows that integrality, D-finiteness, small growth, and three singular
directions alone do not force a Hankel upper rate below $\sqrt e$.
Together with Appendix A, it explains why a proof based on independently
improving the two uniform bounds needs genuinely additional input.

\begin{thebibliography}{9}
\bibitem{Ruzsa} I. Z. Ruzsa, On congruence preserving functions,
\emph{Mat. Lapok} \textbf{22} (1971), 125--134.
\bibitem{Delaygue} E. Delaygue, On Ruzsa's conjecture on congruence preserving
functions, \emph{Bull. London Math. Soc.} \textbf{58} (2026), e70332.
\href{https://arxiv.org/abs/2502.13068}{arXiv:2502.13068v2}.
\bibitem{DR} E. Delaygue and T. Rivoal, On primary pseudo-polynomials
(Around Ruzsa's Conjecture), \emph{Int. J. Number Theory} \textbf{18} (2022),
1613--1636. \href{https://arxiv.org/abs/2102.01534}{arXiv:2102.01534}.
\bibitem{DvK} J. J. Duistermaat and W. van der Kallen,
Constant terms in powers of a Laurent polynomial,
\emph{Indag. Math. (N.S.)} \textbf{9} (1998), 221--231.
\href{https://doi.org/10.1016/S0019-3577(98)80020-7}{doi:10.1016/S0019-3577(98)80020-7}.
\bibitem{BSY} A. Bostan, A. Straub, and S. Yurkevich,
On the representability of sequences as constant terms,
\emph{J. Number Theory} \textbf{253} (2023), 235--256.
\href{https://doi.org/10.1016/j.jnt.2023.06.015}{doi:10.1016/j.jnt.2023.06.015}.
\end{thebibliography}
\end{document}
